\documentclass[12pt]{article}
\usepackage[utf8]{inputenc}
\usepackage[spanish]{babel}
\usepackage{amsmath, amssymb, amsfonts}
\usepackage{geometry}
\geometry{a4paper, margin=2.5cm}

\title{\textbf{Métodos Matemáticos de la Física II} \\ Examen Final – 04/08/2026}
\author{Nombre: \underline{\hspace{6cm}}}
\date{}

\begin{document}
\maketitle

\section*{Problema 1}
Sea \(A\) una matriz \(4 \times 4\) cuyo polinomio característico es
\[
p_A(\lambda) = (\lambda - 2)^2(\lambda + 1)^2.
\]
Además, se sabe que
\[
\operatorname{rang}(A - 2I) = 2, \qquad \operatorname{rang}(A + I) = 3.
\]
\begin{enumerate}
    \item[(a)] Determine las multiplicidades geométricas de cada autovalor.
    \item[(b)] Encuentre la forma canónica de Jordan de \(A\) (no se pide la matriz de cambio de base).
    \item[(c)] ¿Cuántos subespacios invariantes irreducibles tiene la descomposición de \(\mathbb{R}^4\) bajo la acción de \(A\)? Indique la dimensión de cada uno.
\end{enumerate}

\vspace{0.5cm}
{\bf Soluci\'on }{\em 1)}
\vspace{0.5cm}

{\bf (a) Multiplicidades geométricas}

Los autovalores son
$$
\lambda_1=2,\qquad \lambda_2=-1,
$$
ambos con multiplicidad algebraica
$$
m_a(2)=2,\qquad m_a(-1)=2.
$$
La multiplicidad geométrica es
$$
m_g(\lambda)=\dim\ker(A-\lambda I).
$$
Usamos rango-nulidad:
$$
\dim\ker(A-2I)
=
4-\operatorname{rang}(A-2I)
=
4-2
=
2.
$$
Por tanto
$$
m_g(2)=2.
$$
Asimismo,
$$
\dim\ker(A+I)
=
4-\operatorname{rang}(A+I)
=
4-3
=
1,
$$
luego
$$
m_g(-1)=1.
$$

{\em Respuesta:}

$$
\boxed{m_g(2)=2,\qquad m_g(-1)=1.}
$$

{\bf (b) Forma canónica de Jordan}

Para $\lambda=2$, tenemos
$$
m_a(2)=2,\qquad m_g(2)=2.
$$
El número de bloques de Jordan asociados a $\lambda=2$ es igual a $m_g(2)$, es decir, 2 bloques.
Como la suma de sus tamaños debe ser 2, necesariamente son dos bloques $1\times 1$:
$$
J_{2}=
\begin{pmatrix}
2&0\
0&2
\end{pmatrix}.
$$

Para $\lambda=-1$, tenemos
$$
m_a(-1)=2,\qquad m_g(-1)=1.
$$
Hay un único bloque de Jordan asociado a $-1$, y su tamaño debe ser 2.
Por tanto,
$$
J_{-1}
=
\begin{pmatrix}
-1&1\
0&-1
\end{pmatrix}.
$$

{\em Forma de Jordan total.}
Uniendo ambos sectores:
$$
\boxed{
J=
\begin{pmatrix}
2&0&0&0\\
0&2&0&0\\
0&0&-1&1\\
0&0&0&-1
\end{pmatrix}.
}
$$

{\bf (c) Subespacios invariantes irreducibles}

En la forma de Jordan, cada bloque de Jordan corresponde a un subespacio invariante indecomponible (irreducible en el sentido usado en Álgebra Lineal para la descomposición de Jordan).

Tenemos:
\begin{itemize}
\item un bloque $J_1(2)$,
\item otro bloque $J_1(2)$,
\item un bloque $J_2(-1)$.
\end{itemize}
Por lo tanto la descomposición de $\mathbb R^4$ es
$$
\mathbb R^4
=
V_1\oplus V_2\oplus V_3,
$$
con
$$
\dim V_1=1,\qquad
\dim V_2=1,\qquad
\dim V_3=2.
$$
Así, el número de subespacios invariantes irreducibles es
$$
\boxed{3}
$$
y sus dimensiones son
$$
\boxed{1,\;1,\;2.}
$$

{\em Resumen final}

$$
\boxed{m_g(2)=2,\qquad m_g(-1)=1}
$$
$$
\boxed{
J=
\begin{pmatrix}
2&0&0&0\\
0&2&0&0\\
0&0&-1&1\\
0&0&0&-1
\end{pmatrix}
}
$$
$$
\boxed{\mathbb R^4=V_1\oplus V_2\oplus V_3,\qquad
\dim(V_1,V_2,V_3)=(1,1,2).}
$$







\section*{Problema 2}
Considere la superficie de un paraboloide de revolución parametrizada por
\[
x = r \cos \theta, \qquad y = r \sin \theta, \qquad z = \frac{r^2}{2},
\]
con \(r > 0\) y \(0 \le \theta < 2\pi\). La métrica inducida es:
\[
ds^2 = (1 + r^2)\,dr^2 + r^2\,d\theta^2.
\]
\begin{enumerate}
    \item[(a)] Calcule la métrica inversa \(g^{ij}\).
    \item[(b)] Determine todos los símbolos de Christoffel \(\Gamma^k_{ij}\) no nulos.
    \item[(c)] Obtenga el operador Laplaciano \(\Delta u(r, \theta) = g^{ij}\nabla_i\nabla_j u\) para una función escalar \(u(r, \theta)\).
\end{enumerate}

\vspace{0.5cm}
{\bf Soluci\'on }{\em 2)}
\vspace{0.5cm}

El enunciado establece que el paraboloide de revolución está parametrizado por 
$$
x=r\cos\theta,\qquad
y=r\sin\theta,\qquad
z=\frac{r^2}{2},
$$
y que la métrica inducida es
$$
ds^2=(1+r^2),dr^2+r^2,d\theta^2.
$$
Por lo tanto,
$$
(g_{ij})=
\begin{pmatrix}
1+r^2 & 0\\
0 & r^2
\end{pmatrix},
$$
donde las coordenadas son
$$
x^1=r,\qquad x^2=\theta.
$$

{\bf (a) Métrica inversa}

Como la métrica es diagonal, la inversa se obtiene inmediatamente:
$$
(g^{ij})
=
\begin{pmatrix}
\dfrac1{1+r^2} & 0\\
0 & \dfrac1{r^2}
\end{pmatrix}.
$$
Por consiguiente,
$$
\boxed{
g^{rr}=\frac1{1+r^2},
\qquad
g^{\theta\theta}=\frac1{r^2},
\qquad
g^{r\theta}=g^{\theta r}=0.
}
$$

{\bf (b) Símbolos de Christoffel}

Usamos
$$
\Gamma^k_{ij}
=
\frac12\,g^{k\ell}
\left(
\partial_i g_{j\ell}
+
\partial_j g_{i\ell}
-
\partial_\ell g_{ij}
\right).
$$
Las únicas derivadas no nulas de la métrica son
$$
\partial_r g_{rr}=2r,
\qquad
\partial_r g_{\theta\theta}=2r.
$$
Todas las derivadas respecto de $\theta$ son cero.

{\em Cálculo de $\Gamma^r_{rr}$}

$$
\Gamma^r_{rr}
=
\frac12 g^{rr}\partial_r g_{rr}
=
\frac12\frac1{1+r^2}(2r)
=
\frac{r}{1+r^2}.
$$

$$
\boxed{\Gamma^r_{rr}=\frac{r}{1+r^2}}
$$

{\em Cálculo de $\Gamma^r_{\theta\theta}$}

$$
\Gamma^r_{\theta\theta}
=
-\frac12 g^{rr}\partial_r g_{\theta\theta}
=
-\frac12\frac1{1+r^2}(2r)
=
-\frac{r}{1+r^2}.
$$

$$
\boxed{\Gamma^r_{\theta\theta}
=
-\frac{r}{1+r^2}}
$$

{\em Cálculo de $\Gamma^\theta_{r\theta}$}

$$
\Gamma^\theta_{r\theta}
=
\Gamma^\theta_{\theta r}
=
\frac12 g^{\theta\theta}
\partial_r g_{\theta\theta}
=
\frac12\frac1{r^2}(2r)
=
\frac1r.
$$

$$
\boxed{
\Gamma^\theta_{r\theta}
=
\Gamma^\theta_{\theta r}
=
\frac1r
}
$$

{\em Los demás}

Todos los demás símbolos se anulan.
En resumen,
$$
\boxed{
\Gamma^r_{rr}=\frac{r}{1+r^2},
\qquad
\Gamma^r_{\theta\theta}
=
-\frac{r}{1+r^2},
\qquad
\Gamma^\theta_{r\theta}
=
\Gamma^\theta_{\theta r}
=
\frac1r.
}
$$

{\bf (c) Laplaciano}

Para una función escalar $u(r,\theta)$,
$$
\nabla_i\nabla_j u
=
\partial_i\partial_j u
-
\Gamma^k_{ij}\partial_k u.
$$
El Laplaciano es
$$
\Delta u
=
g^{ij}\nabla_i\nabla_j u.
$$
Como la métrica es diagonal,
$$
\Delta u
=
g^{rr}\nabla_r\nabla_r u
+
g^{\theta\theta}\nabla_\theta\nabla_\theta u.
$$

{\em Parte radial}

$$
\nabla_r\nabla_r u
=
u_{rr}
=
\Gamma^r_{rr}u_r
-
u_{rr}
=
\frac{r}{1+r^2}u_r.
$$
Por tanto
$$
g^{rr}\nabla_r\nabla_r u
=
\frac{u_{rr}}{1+r^2}
-
\frac{r}{(1+r^2)^2}u_r.
$$

{\em Parte angular}

$$
\nabla_\theta\nabla_\theta u
=
u_{\theta\theta}
-
\Gamma^r_{\theta\theta}u_r.
$$
Como
$$
\Gamma^r_{\theta\theta}
=
-\frac{r}{1+r^2},
$$
resulta
$$
\nabla_\theta\nabla_\theta u
=
u_{\theta\theta}
+
\frac{r}{1+r^2}u_r.
$$
Multiplicando por $g^{\theta\theta}=1/r^2$,
$$
g^{\theta\theta}\nabla_\theta\nabla_\theta u
=
\frac{u_{\theta\theta}}{r^2}
+
\frac1{r(1+r^2)}u_r.
$$

{\em Sumando}

$$
\Delta u
=
\frac{u_{rr}}{1+r^2}
+
\frac{u_{\theta\theta}}{r^2}
+
\left[
\frac1{r(1+r^2)}
-
\frac{r}{(1+r^2)^2}
\right]u_r.
$$
Llevando al mismo denominador:
$$
\frac1{r(1+r^2)}
-
\frac{r}{(1+r^2)^2}
=
\frac1{r(1+r^2)^2}.
$$
Finalmente,
$$
\boxed{
\Delta u(r,\theta)
=
\frac1{1+r^2}\,
\frac{\partial^2u}{\partial r^2}
+
\frac1{r^2}\,
\frac{\partial^2u}{\partial\theta^2}
+
\frac1{r(1+r^2)^2}\,
\frac{\partial u}{\partial r}
}.
$$
Como verificación, si se utiliza la fórmula alternativa
$$
\Delta u
=
\frac1{\sqrt g}
\partial_i
\left(
\sqrt g,g^{ij}\partial_j u
\right),
\qquad
\sqrt g=r\sqrt{1+r^2},
$$
se obtiene exactamente el mismo resultado.

{\bf Respuesta final:}

$$
g^{ij}
=
\begin{pmatrix}
\dfrac1{1+r^2} & 0\\
0 & \dfrac1{r^2}
\end{pmatrix},
$$

$$
\Gamma^r_{rr}=\frac{r}{1+r^2},
\qquad
\Gamma^r_{\theta\theta}=-\frac{r}{1+r^2},
\qquad
\Gamma^\theta_{r\theta}
=
\Gamma^\theta_{\theta r}
=
\frac1r,
$$

$$
\boxed{
\Delta u
=
\frac{u_{rr}}{1+r^2}
+
\frac{u_{\theta\theta}}{r^2}
+
\frac{1}{r(1+r^2)^2}\,u_r.
}
$$








\section*{Problema 3}
Utilice el método de las características para resolver el siguiente problema de Cauchy:
\[
\frac{\partial u}{\partial t} + 2t \frac{\partial u}{\partial x} = u, \qquad u(x, 0) = \sin x.
\]

\vspace{0.5cm}
{\bf Soluci\'on }{\em 3)}
\vspace{0.5cm}

{\bf 1. Ecuaciones características}

La ecuación es de la forma
$$
u_t+a(x,t)u_x=b(x,t,u),
$$
con
$$
a(x,t)=2t,
\qquad
b(x,t,u)=u.
$$
Las ecuaciones características son
$$
\frac{dt}{ds}=1,
$$
$$
\frac{dx}{ds}=2t,
$$
$$
\frac{du}{ds}=u.
$$
Además, la curva inicial es
$$
t=0,
\qquad
x=\alpha,
\qquad
u=\sin\alpha,
$$

donde $\alpha$ parametriza los puntos de la condición inicial.

{\bf 2. Resolver la ecuación para $t$}

De
$$
\frac{dt}{ds}=1,
$$
obtenemos
$$
t=s+C.
$$
Imponiendo que sobre la curva inicial (t=0) cuando $s=0$,
$$
C=0.
$$
Por lo tanto,
$$
\boxed{t=s}.
$$

{\bf 3. Resolver la ecuación para $x$}

Sustituyendo $t=s$,
$$
\frac{dx}{ds}=2s.
$$
Integrando,
$$
x=s^2+C_1.
$$
Sobre la curva inicial,
$$
x(0)=\alpha,
$$
de modo que
$$
C_1=\alpha.
$$
Así,
$$
\boxed{x=s^2+\alpha}.
$$
Como $s=t$,
$$
\boxed{x=t^2+\alpha}.
$$
Luego
$$
\boxed{\alpha=x-t^2}.
$$

{\bf 4. Resolver la ecuación para $u$}

Tenemos
$$
\frac{du}{ds}=u.
$$
Integrando,
$$
u=C_2e^s.
$$
Sobre la curva inicial,
$$
u(0)=\sin\alpha,
$$
por lo que
$$
C_2=\sin\alpha.
$$
Entonces
$$
u(s)=e^s\sin\alpha.
$$
Como $s=t$,
$$
u=e^t\sin\alpha.
$$
Finalmente sustituimos
$$
\alpha=x-t^2.
$$

{\bf 5. Solución}

$$
\boxed{
u(x,t)
=
e^t\sin(x-t^2).
}
$$

{\bf Verificación}

Calculemos:
$$
u_x=e^t\cos(x-t^2),
$$
$$
u_t
=
e^t\sin(x-t^2)
-2t,e^t\cos(x-t^2).
$$
Entonces
$$
u_t+2t\,u_x
=
e^t\sin(x-t^2)
=u.
$$
Además,
$$
u(x,0)=e^0\sin(x)=\sin x.
$$
La solución satisface tanto la EDP como la condición inicial.

{\bf Respuesta final}

$$
\boxed{
u(x,t)=e^t\sin(x-t^2).
}
$$
El punto clave es identificar correctamente la constante característica $\alpha=x-t^2$.







\section*{Problema 4}
En el interior de un cilindro infinito de radio \(R\), el potencial \(\Phi(r, \theta)\) satisface la ecuación de Laplace. En la pared \(r = R\) se impone la condición
\[
\Phi(R, \theta) = V_0 \cos^2 \theta,
\]
donde \(V_0\) es una constante.
\begin{enumerate}
    \item[(a)] Exprese \(\cos^2 \theta\) como combinación lineal de armónicos circulares (es decir, en términos de senos y cosenos de \(n\theta\)).
    \item[(b)] Encuentre \(\Phi(r, \theta)\) en el interior del cilindro, como una serie de potencias de \(r\) (finita, en este caso).
    \item[(c)] Escriba la solución en forma cerrada.
\end{enumerate}

\vspace{0.5cm}
{\bf Soluci\'on }{\em 4)}
\vspace{0.5cm}

{\bf (a) Expresar $\cos^2\theta$ en armónicos circulares}

Utilizamos la identidad trigonométrica
$$
\cos^2\theta
=
\frac{1+\cos(2\theta)}{2}.
$$
Por tanto,
$$
\Phi(R,\theta)
=
\frac{V_0}{2}
+
\frac{V_0}{2}\cos(2\theta).
$$
La condición de borde contiene únicamente los modos
$$
n=0,
\qquad
n=2.
$$

{\bf (b) Solución general de Laplace en el disco}

La ecuación de Laplace en coordenadas polares es
$$
\nabla^2\Phi
=
\frac{\partial^2\Phi}{\partial r^2}
+\frac1r\frac{\partial\Phi}{\partial r}
+\frac1{r^2}
\frac{\partial^2\Phi}{\partial\theta^2}
=0.
$$
La solución regular en el origen tiene la forma
$$
\Phi(r,\theta)
=
A_0
+
\sum_{n=1}^{\infty}
r^n
\left(
A_n\cos n\theta
+
B_n\sin n\theta
\right).
$$
Los términos $r^{-n}$ se descartan porque divergen en $r=0$.
Como el borde sólo contiene un término constante y un $\cos(2\theta)$, la solución debe ser
$$
\Phi(r,\theta)
=
A_0+A_2r^2\cos(2\theta).
$$

{\bf Determinación de los coeficientes}

Imponemos la condición en $r=R$:
$$
A_0+A_2R^2\cos(2\theta)
=
\frac{V_0}{2}
+
\frac{V_0}{2}\cos(2\theta).
$$
Comparando coeficientes:

{\em Modo constante}

$$
A_0=\frac{V_0}{2}.
$$

{\em Modo $n=2$}

$$
A_2R^2=\frac{V_0}{2},
$$
de donde
$$
A_2=\frac{V_0}{2R^2}.
$$

{\bf (c) Forma cerrada}

Sustituyendo:
$$
\boxed{
\Phi(r,\theta)
=
\frac{V_0}{2}
+
\frac{V_0}{2R^2}r^2\cos(2\theta).
}
$$
También puede escribirse como
$$
\boxed{
\Phi(r,\theta)
=
\frac{V_0}{2}
\left(
1+\frac{r^2}{R^2}\cos(2\theta)
\right).
}
$$

{\bf Verificación}

En el borde $r=R$:
$$
\Phi(R,\theta)
=
\frac{V_0}{2}
\left(1+\cos2\theta\right)
=
V_0\cos^2\theta,
$$
usando nuevamente
$$
\cos^2\theta=\frac{1+\cos2\theta}{2}.
$$
Por lo tanto la solución satisface exactamente la condición de contorno.

{\bf Respuesta final}

$$
\boxed{
\Phi(r,\theta)
=
\frac{V_0}{2}
\left(
1+\frac{r^2}{R^2}\cos 2\theta
\right).
}
$$

Observación: este es uno de los ejercicios más sencillos de separación de variables en el disco porque la condición de borde contiene un único armónico no trivial $(n=2)$, por lo que la serie de Fourier se trunca inmediatamente.








\section*{Problema 5}
Resuelva la ecuación de Poisson \(\nabla^2 u(r, \phi) = \beta_0 \sen(\phi)\) (con \(\beta_0\) constante) dentro de un disco de radio \(R\), sabiendo que en su borde la condición es \(u(R, \phi) = \cos(\phi)\).

\vspace{0.5cm}
{\bf Soluci\'on }{\em 5)}
\vspace{0.5cm}

{\bf 1. Idea general}

La ecuación es de Poisson, así que escribimos
$$
u=u_p+u_h,
$$
donde
$$
\nabla^2u_p=\beta_0\sin\varphi
$$
es una solución particular, y
$$
\nabla^2u_h=0
$$
es una solución armónica que ajustará la condición de borde.

{\bf 2. Buscar una solución particular}

El término fuente tiene dependencia angular $\sin\varphi$, por lo que intentamos
$$
u_p(r,\varphi)=f(r)\sin\varphi.
$$
En coordenadas polares,
$$
\nabla^2u=
u_{rr}
+\frac1r u_r
+\frac1{r^2}u_{\varphi\varphi}.
$$
Calculamos:
$$
u_r=f'(r)\sin\varphi,
$$
$$
u_{rr}=f''(r)\sin\varphi,
$$
$$
u_{\varphi\varphi}
=-f(r)\sin\varphi.
$$
Entonces
$$
\nabla^2u_p
=
\left(
f''+\frac1r f'
-\frac{f}{r^2}
\right)\sin\varphi.
$$
Queremos que esto sea
$$
\beta_0\sin\varphi,
$$
por lo que $f$ debe satisfacer
$$
f''+\frac1r f'
-\frac{f}{r^2}
=
\beta_0.
$$

{\bf 3. Resolver para $f(r)$}

Probamos un polinomio
$$
f(r)=Ar^2.
$$
Entonces
$$
f'=2Ar,
\qquad
f''=2A.
$$
Sustituyendo:
$$
2A+\frac1r(2Ar)-\frac{Ar^2}{r^2}
=
2A+2A-A
=
3A.
$$
Para obtener $\beta_0$,
$$
3A=\beta_0
\quad\Longrightarrow\quad
A=\frac{\beta_0}{3}.
$$
Por tanto,
$$
\boxed{
u_p(r,\varphi)
=
\frac{\beta_0}{3}r^2\sin\varphi.
}
$$

{\bf 4. Solución de Laplace}

La solución regular de Laplace en el disco es
$$
u_h
=
A_0
+
\sum_{n=1}^{\infty}
r^n
\left(
A_n\cos n\varphi
+
B_n\sin n\varphi
\right).
$$
La condición de borde contiene solamente un modo $\cos\varphi$, por lo que bastará tomar
$$
u_h
=
A\,r\cos\varphi
+
B\,r\sin\varphi .
$$

{\bf 5. Imponer la condición de borde}

La solución total es
$$
u(r,\varphi)
=
A r\cos\varphi
+
B r\sin\varphi
+
\frac{\beta_0}{3}r^2\sin\varphi.
$$
En $r=R$,
$$
A R\cos\varphi
+
\left(
BR+\frac{\beta_0}{3}R^2
\right)\sin\varphi
=
\cos\varphi.
$$
Igualando coeficientes de Fourier:

{\em Modo $\cos\varphi$}

$$
AR=1
\quad\Longrightarrow\quad
A=\frac1R.
$$

{\bf Modo $\sin\varphi$}

$$
BR+\frac{\beta_0}{3}R^2=0,
$$
de donde
$$
B=-\frac{\beta_0}{3}R.
$$

{\bf 6. Solución final}

Sustituyendo $A$ y $B$,
$$
u(r,\varphi)
=
\frac{r}{R}\cos\varphi
-\frac{\beta_0R}{3}\,r\sin\varphi
+\frac{\beta_0}{3}r^2\sin\varphi.
$$
Agrupando los términos con $\sin\varphi$,
$$
\boxed{
u(r,\varphi)
=
\frac{r}{R}\cos\varphi
+
\frac{\beta_0}{3}
\left(r^2-Rr\right)\sin\varphi.
}
$$

{\bf Verificación}

En el borde:
$$
u(R,\varphi)
=
\cos\varphi
+
\frac{\beta_0}{3}(R^2-R^2)\sin\varphi
=
\cos\varphi.
$$
Además,
$$
\nabla^2\left(\frac{r}{R}\cos\varphi\right)=0,
$$
$$
\nabla^2\left[\frac{\beta_0}{3}(r^2-Rr)\sin\varphi\right]
=
\beta_0\sin\varphi,
$$
porque $r\sin\varphi$ es armónica.
Por lo tanto,
$$
\boxed{
u(r,\varphi)
=
\frac{r}{R}\cos\varphi
+
\frac{\beta_0}{3}(r^2-Rr)\sin\varphi
}
$$
es la solución buscada.







\section*{Problema 6}

Encuentre la función de Green correspondiente a la ecuación diferencial ordinaria
\[
x^2 y'' + 4xy' + (x^2 + 2)y = 0, \qquad 1 < x < 2,
\]
con condiciones de contorno
\[
y(1) = 0, \qquad y(2) = 0.
\]
\textit{Ayuda:} Considere el cambio de variable \(v(x) = x^2 y(x)\).

\vspace{0.5cm}
{\bf Soluci\'on }{\em 6)}
\vspace{0.5cm}

{\bf 1. Transformación de la ecuación}

Escribimos
$$
y=\frac{v}{x^2}.
$$
Entonces
$$
y'
=\frac{v'}{x^2}-\frac{2v}{x^3},
$$
$$
y''
=\frac{v''}{x^2}
-\frac{4v'}{x^3}
+\frac{6v}{x^4}.
$$
Sustituyendo en la ecuación:
$$
x^2y''
=
v''
-\frac{4}{x}v'
+\frac{6}{x^2}v,
$$
$$
4xy'
=
\frac{4}{x}v'
-\frac{8}{x^2}v,
$$
$$
(x^2+2)y
=
v+\frac{2}{x^2}v.
$$
Sumando:
$$
v''
-\frac{4}{x}v'
+\frac{6}{x^2}v
+\frac{4}{x}v'
-\frac{8}{x^2}v
+v+\frac{2}{x^2}v.
$$
Los términos con $v'$ se cancelan y también los de $1/x^2$:
$$
v''+v=0.
$$
Así, el operador original queda reducido a
$$
\boxed{v''+v=0.}
$$

{\bf 2. Condiciones de contorno}

Como
$$
v=x^2y,
$$
resulta
$$
v(1)=1^2y(1)=0,
$$
$$
v(2)=2^2y(2)=0.
$$
Por tanto:
$$
\boxed{
v''+v=0,
\qquad
v(1)=v(2)=0.
}
$$

{\bf 3. Soluciones homogéneas}

Dos soluciones independientes son
$$
u_1(x)=\sin x,
\qquad
u_2(x)=\cos x.
$$
Necesitamos soluciones adaptadas a cada borde.

{\em Solución izquierda}

Que satisfaga $v(1)=0$:
$$
\phi(x)=\sin(x-1).
$$

{\em Solución derecha}

Que satisfaga $v(2)=0$:
$$
\psi(x)=\sin(2-x).
$$

{\bf 4. Wronskiano}

Calculamos
$$
W=\phi\psi'-\phi'\psi.
$$
Como
$$
\phi'=\cos(x-1),
$$
$$
\psi'=-\cos(2-x),
$$
obtenemos
$$
W
=
-\sin(x-1)\cos(2-x)
-\cos(x-1)\sin(2-x).
$$
Usando
$$
\sin A\cos B+\cos A\sin B
=
\sin(A+B),
$$
obtenemos que el Wronskiano
$$
W
=
-\sin[(x-1)+(2-x)]
=
-\sin 1.
$$
resulta constante.

{\bf 5. Función de Green para $v''+v$}

La fórmula estándar es
$$
G_v(x,\xi)
=
\frac{\phi(x_<)\psi(x_>)}{W},
$$
donde
$$
x_< = \min(x,\xi),
\qquad
x_> = \max(x,\xi).
$$
Luego
$$
\boxed{
G_v(x,\xi)
=
-\frac{
\sin(x_<-1),
\sin(2-x_>)
}{
\sin 1
}.
}
$$
Equivalentemente,
$$
G_v(x,\xi)=
\begin{cases}
-\dfrac{\sin(x-1)\sin(2-\xi)}
{\sin1},
& x<\xi,
\\
-\dfrac{\sin(\xi-1)\sin(2-x)}
{\sin1},
& x>\xi.
\end{cases}
$$

{\bf 6. Volver a la variable $y$}

La ecuación original puede escribirse como
$$
L_y[y]=f(x).
$$
Multiplicando por $x^2$,
$$
v''+v=x^2f(x).
$$
Por definición,
$$
v(x)
=
\int_1^2
G_v(x,\xi),
\xi^2f(\xi),
d\xi.
$$
Como
$$
y(x)=\frac{v(x)}{x^2},
$$
la función de Green del operador original es
$$
G_y(x,\xi)
=
\frac{\xi^2}{x^2}
,G_v(x,\xi).
$$
Por tanto,
$$
\boxed{
G(x,\xi)
=
-\frac{\xi^2}{x^2\sin1}
\begin{cases}
\sin(x-1)\sin(2-\xi),
&
x<\xi,
\\
\sin(\xi-1)\sin(2-x),
&
x>\xi.
\end{cases}
}
$$

{\bf Respuesta final}

La función de Green asociada al problema de contorno
$$
x^2y''+4xy'+(x^2+2)y=f(x),
\qquad
y(1)=y(2)=0,
$$
es
$$
\boxed{
G(x,\xi)
=
-\frac{\xi^2}{x^2\sin1}
\begin{cases}
\sin(x-1)\sin(2-\xi), & x<\xi,\\
\sin(\xi-1)\sin(2-x), & x>\xi.
\end{cases}
}
$$
La clave del ejercicio era reconocer que el cambio sugerido $v=x^2y$ elimina completamente el término de primera derivada y transforma la ecuación en el problema elemental $v''+v=0$.









\section*{Problema 7}
Una cuerda elástica de longitud \(L\) se halla fija en sus extremos y sumergida en un fluido que presenta resistencia a su movimiento, de modo que la desviación \(u(x, t)\) satisface la ecuación
\[
u_{tt} + \gamma u_t - c^2 u_{xx} = 0,
\]
con \(\gamma\) una constante positiva y condiciones de contorno \(u(0,t) = u(L,t) = 0\).
\begin{enumerate}
    \item[(a)] Proponiendo \(u(x,t) = X(x)T(t)\), escriba las ecuaciones diferenciales satisfechas por \(X\) y por \(T\).
    \item[(b)] Muestre que las autofunciones espaciales que satisfacen las condiciones de contorno son de la forma \(X_n(x) = \sen(k_n x)\), y determine \(k_n\).
    \item[(c)] Proponiendo ahora una solución en la forma
    \[
    u(x, t) = \sum_{n=1}^{\infty} A_n X_n(x) T_n(t),
    \]
    determine la ecuación diferencial satisfecha por cada función \(T_n(t)\).
    \item[(d)] Muestre que \(T_n(t)\) oscila (amortiguadamente) sólo si \(n < \dfrac{L\gamma}{\pi c}\).
\end{enumerate}

\vspace{0.5cm}
{\bf Soluci\'on }{\em 7)}
\vspace{0.5cm}

La ecuación de movimiento de la cuerda es

\begin{equation}
\frac{\partial^2u}{\partial t^2}
+\gamma\frac{\partial u}{\partial t}
-c^2\frac{\partial^2u}{\partial x^2}=0,
\qquad
0<x<L,
\end{equation}

con condiciones de contorno

\begin{equation}
u(0,t)=u(L,t)=0,
\end{equation}

donde $\gamma>0$ representa el coeficiente de amortiguamiento del fluido.

\subsubsection*{(a) Separación de variables}

Buscamos una solución de la forma

\begin{equation}
u(x,t)=X(x)T(t).
\end{equation}

Entonces,

\[
u_t=XT', \qquad
u_{tt}=XT'', \qquad
u_{xx}=X''T.
\]

Sustituyendo en la ecuación diferencial,

\begin{equation}
XT''+\gamma XT'
-c^2X''T=0.
\end{equation}

Dividiendo por $XT$,

\begin{equation}
\frac{T''}{T}
+\gamma\frac{T'}{T}
=
c^2\frac{X''}{X}.
\end{equation}

El miembro izquierdo depende únicamente de $t$, mientras que el derecho depende únicamente de $x$. Ambos deben ser iguales a una constante de separación.

Elegimos la constante como

\[
-c^2k^2,
\]

de modo que el problema espacial resulte oscilatorio.

Obtenemos las ecuaciones

\begin{equation}
X''+k^2X=0,
\end{equation}

y

\begin{equation}
T''+\gamma T'+c^2k^2T=0.
\end{equation}

Éstas son las ecuaciones pedidas.

\subsubsection*{(b) Autofunciones espaciales}

Debemos resolver

\begin{equation}
X''+k^2X=0,
\end{equation}

con condiciones de contorno

\begin{equation}
X(0)=0,
\qquad
X(L)=0.
\end{equation}

La solución general es

\begin{equation}
X(x)=A\cos(kx)+B\sin(kx).
\end{equation}

La primera condición implica

\[
A=0,
\]

por lo que

\begin{equation}
X(x)=B\sin(kx).
\end{equation}

Aplicando la segunda condición,

\[
B\sin(kL)=0.
\]

Para obtener soluciones no triviales ($B\neq0$), debe cumplirse

\begin{equation}
\sin(kL)=0,
\end{equation}

de donde

\begin{equation}
kL=n\pi,
\qquad
n=1,2,3,\ldots
\end{equation}

Por consiguiente,

\begin{equation}
k_n=\frac{n\pi}{L},
\end{equation}

y las autofunciones espaciales son

\begin{equation}
X_n(x)=
\sin\!\left(\frac{n\pi x}{L}\right).
\end{equation}

Estas funciones forman una base ortogonal en el intervalo $(0,L)$.

\subsubsection*{(c) Desarrollo modal}

Proponemos la solución

\begin{equation}
u(x,t)=
\sum_{n=1}^{\infty}
A_nX_n(x)T_n(t).
\end{equation}

Como

\[
X_n''=-k_n^2X_n,
\]

cada modo satisface independientemente la ecuación

\begin{equation}
T_n''+\gamma T_n'
+c^2k_n^2T_n=0.
\end{equation}

Sustituyendo el valor de $k_n$,

\begin{equation}
T_n''+\gamma T_n'
+\left(\frac{n\pi c}{L}\right)^2
T_n=0.
\end{equation}

Ésta es la ecuación diferencial buscada.

\subsubsection*{(d) Condición para que exista movimiento oscilatorio}

Cada modo satisface la ecuación

\begin{equation}
T_n''+\gamma T_n'
+\omega_n^2T_n=0,
\end{equation}

donde

\begin{equation}
\omega_n=\frac{n\pi c}{L}.
\end{equation}

La ecuación característica es

\begin{equation}
r^2+\gamma r+\omega_n^2=0,
\end{equation}

cuyas raíces son

\begin{equation}
r=
-\frac{\gamma}{2}
\pm
\frac12
\sqrt{\gamma^2-4\omega_n^2}.
\end{equation}

Existen tres posibilidades.

\paragraph{Sobreamortiguamiento.}

Si

\[
\gamma^2>4\omega_n^2,
\]

las raíces son reales y negativas. En consecuencia, no existen oscilaciones.

\paragraph{Amortiguamiento crítico.}

Si

\[
\gamma^2=4\omega_n^2,
\]

aparece una raíz doble,

\[
r=-\frac{\gamma}{2},
\]

y tampoco existen oscilaciones.

\paragraph{Subamortiguamiento.}

Si

\[
\gamma^2<4\omega_n^2,
\]

las raíces son complejas,

\begin{equation}
r=
-\frac{\gamma}{2}
\pm
i
\sqrt{\omega_n^2-\frac{\gamma^2}{4}}.
\end{equation}

La solución temporal es

\begin{equation}
T_n(t)=
e^{-\gamma t/2}
\left[
C_n\cos(\Omega_nt)
+
D_n\sin(\Omega_nt)
\right],
\end{equation}

donde

\begin{equation}
\Omega_n=
\sqrt{
\omega_n^2-\frac{\gamma^2}{4}
}.
\end{equation}

Por lo tanto, el modo oscila amortiguadamente únicamente cuando

\begin{equation}
\omega_n>\frac{\gamma}{2}.
\end{equation}

Como

\[
\omega_n=\frac{n\pi c}{L},
\]

se obtiene

\begin{equation}
\frac{n\pi c}{L}>
\frac{\gamma}{2},
\end{equation}

es decir,

\begin{equation}
\boxed{
n>
\frac{\gamma L}{2\pi c}.
}
\end{equation}

\paragraph{Observación.}

El enunciado solicita demostrar que

\[
n<
\frac{\gamma L}{\pi c}.
\]

Sin embargo, el análisis de la ecuación característica conduce inequívocamente a la condición

\[
n>
\frac{\gamma L}{2\pi c}.
\]

Por ello, la desigualdad indicada en el enunciado parece contener un error tipográfico, tanto en el sentido de la desigualdad como en el factor numérico.









\section*{Problema 8}
Encuentre la temperatura \(u(r, \phi, t)\) en un cuarto de círculo de radio \(a\) que satisface la ecuación de calor:
\[
\frac{\partial u}{\partial t} = k \nabla^2 u,
\]
sujeto a las condiciones de frontera:
\[
u(r, 0, t) = 0, \qquad u(a, \phi, t) = 0, \qquad u(r, \pi/2, t) = 0,
\]
y a la condición inicial:
\[
u(r, \phi, 0) = r \sen(4\phi).
\]

\vspace{0.5cm}
{\bf Soluci\'on }{\em 8)}
\vspace{0.5cm}

Debemos resolver la ecuación del calor

\begin{equation}
\frac{\partial u}{\partial t}=k\nabla^2u,
\end{equation}

en un cuarto de disco de radio $a$, con condiciones de contorno

\begin{equation}
u(r,0,t)=0,\qquad
u(r,\pi/2,t)=0,\qquad
u(a,\phi,t)=0,
\end{equation}

y condición inicial

\begin{equation}
u(r,\phi,0)=r\sin(4\phi).
\end{equation}

Recordemos que el Laplaciano en coordenadas polares es

\begin{equation}
\nabla^2u=
\frac{\partial^2u}{\partial r^2}
+\frac1r\frac{\partial u}{\partial r}
+\frac1{r^2}\frac{\partial^2u}{\partial\phi^2}.
\end{equation}

\subsubsection*{Separación de variables}

Buscamos una solución de la forma

\begin{equation}
u(r,\phi,t)=R(r)\Phi(\phi)T(t).
\end{equation}

Sustituyendo en la ecuación del calor,

\[
R\Phi T'
=
k
\left(
R''\Phi T
+\frac1rR'\Phi T
+\frac1{r^2}R\Phi''T
\right).
\]

Dividiendo por $kR\Phi T$,

\begin{equation}
\frac1k\frac{T'}{T}
=
\frac{R''}{R}
+\frac1r\frac{R'}{R}
+\frac1{r^2}\frac{\Phi''}{\Phi}.
\end{equation}

Multiplicando por $r^2$,

\begin{equation}
\frac{r^2}{k}\frac{T'}{T}
=
r^2\frac{R''}{R}
+r\frac{R'}{R}
+\frac{\Phi''}{\Phi}.
\end{equation}

Introducimos una constante de separación $-\lambda$,

\begin{equation}
\frac{\Phi''}{\Phi}=-\lambda,
\end{equation}

obteniéndose

\begin{equation}
\Phi''+\lambda\Phi=0,
\end{equation}

y

\begin{equation}
\frac1k\frac{T'}{T}
=
\frac{R''}{R}
+\frac1r\frac{R'}{R}
-\frac{\lambda}{r^2}.
\end{equation}

Separando nuevamente mediante una constante $-\mu^2$,

\begin{equation}
\frac1k\frac{T'}{T}
=
-\mu^2,
\end{equation}

resultan las ecuaciones

\begin{equation}
T'+k\mu^2T=0,
\end{equation}

y

\begin{equation}
R''+\frac1rR'
+\left(
\mu^2-\frac{\lambda}{r^2}
\right)R=0.
\end{equation}

Esta última es la ecuación de Bessel.

\subsubsection*{Ecuación angular}

Las condiciones de borde son

\[
\Phi(0)=0,
\qquad
\Phi\!\left(\frac{\pi}{2}\right)=0.
\]

La solución general es

\[
\Phi(\phi)=
A\cos(\sqrt{\lambda}\,\phi)
+
B\sin(\sqrt{\lambda}\,\phi).
\]

La condición en $\phi=0$ implica

\[
A=0.
\]

Luego

\[
\Phi(\phi)=
B\sin(\sqrt{\lambda}\,\phi).
\]

La segunda condición exige

\[
\sin\!\left(
\sqrt{\lambda}\,
\frac{\pi}{2}
\right)=0,
\]

de donde

\begin{equation}
\sqrt{\lambda}=2n,
\qquad
n=1,2,\ldots
\end{equation}

Por lo tanto,

\begin{equation}
\lambda_n=4n^2,
\end{equation}

y las autofunciones angulares son

\begin{equation}
\Phi_n(\phi)=\sin(2n\phi).
\end{equation}

\subsubsection*{Ecuación radial}

La ecuación radial queda

\begin{equation}
R''+\frac1rR'
+\left(
\mu^2-\frac{4n^2}{r^2}
\right)R=0.
\end{equation}

Ésta es la ecuación de Bessel de orden $2n$.

Su solución general es

\begin{equation}
R(r)=
AJ_{2n}(\mu r)
+
BY_{2n}(\mu r),
\end{equation}

donde $J_m$ y $Y_m$ son las funciones de Bessel de primera y segunda especie.

Como la solución debe permanecer finita en el origen,

\[
Y_{2n}(r)\rightarrow\infty
\qquad
(r\rightarrow0),
\]

debemos tomar

\[
B=0.
\]

Así,

\begin{equation}
R(r)=J_{2n}(\mu r).
\end{equation}

La condición de borde

\[
u(a,\phi,t)=0
\]

implica

\begin{equation}
J_{2n}(\mu a)=0.
\end{equation}

Si denotamos por

\[
j_{2n,m},
\qquad
m=1,2,\ldots
\]

el $m$-ésimo cero positivo de $J_{2n}$, entonces

\begin{equation}
\mu_{nm}
=
\frac{j_{2n,m}}{a}.
\end{equation}

Las autofunciones radiales son

\begin{equation}
R_{nm}(r)
=
J_{2n}
\!\left(
\frac{j_{2n,m}}{a}r
\right).
\end{equation}

\subsubsection*{Parte temporal}

La ecuación temporal es

\begin{equation}
T'+k\mu_{nm}^2T=0,
\end{equation}

cuya solución es

\begin{equation}
T_{nm}(t)
=
\exp
\!\left[
-k
\left(
\frac{j_{2n,m}}{a}
\right)^2
t
\right].
\end{equation}

\subsubsection*{Solución general}

Combinando todas las soluciones obtenidas,

\begin{equation}
u(r,\phi,t)
=
\sum_{n=1}^{\infty}
\sum_{m=1}^{\infty}
A_{nm}
J_{2n}
\!\left(
\frac{j_{2n,m}}{a}r
\right)
\sin(2n\phi)
\exp
\!\left[
-k
\left(
\frac{j_{2n,m}}{a}
\right)^2
t
\right].
\end{equation}

\subsubsection*{Determinación de los coeficientes}

La condición inicial es

\begin{equation}
u(r,\phi,0)=r\sin(4\phi).
\end{equation}

Como

\[
\sin(4\phi)=\sin(2\cdot2\,\phi),
\]

únicamente contribuye el modo angular correspondiente a

\[
n=2.
\]

En consecuencia,

\begin{equation}
u(r,\phi,t)
=
\sin(4\phi)
\sum_{m=1}^{\infty}
A_m
J_4
\!\left(
\frac{j_{4,m}}{a}r
\right)
\exp
\!\left[
-k
\left(
\frac{j_{4,m}}{a}
\right)^2
t
\right],
\end{equation}

donde los coeficientes se obtienen expandiendo la función radial $r$ en la base de funciones de Bessel,

\begin{equation}
r
=
\sum_{m=1}^{\infty}
A_m
J_4
\!\left(
\frac{j_{4,m}}{a}r
\right).
\end{equation}

Usando la ortogonalidad de las funciones de Bessel,

\begin{equation}
\int_0^a
r
J_4
\!\left(
\frac{j_{4,m}}{a}r
\right)
J_4
\!\left(
\frac{j_{4,n}}{a}r
\right)
dr
=
\frac{a^2}{2}
J_5^2(j_{4,m})
\delta_{mn},
\end{equation}

se obtiene

\begin{equation}
A_m=
\frac{2}
{a^2J_5^2(j_{4,m})}
\int_0^a
r^2
J_4
\!\left(
\frac{j_{4,m}}{a}r
\right)
dr.
\end{equation}

\subsubsection*{Respuesta final}

La temperatura en el cuarto de círculo es

\begin{equation}
\boxed{
u(r,\phi,t)
=
\sin(4\phi)
\sum_{m=1}^{\infty}
A_m
J_4
\!\left(
\frac{j_{4,m}}{a}r
\right)
\exp
\!\left[
-k
\left(
\frac{j_{4,m}}{a}
\right)^2
t
\right],
}
\end{equation}

donde

\begin{equation}
A_m=
\frac{2}
{a^2J_5^2(j_{4,m})}
\int_0^a
r^2
J_4
\!\left(
\frac{j_{4,m}}{a}r
\right)
dr.
\end{equation}

Los autovalores vienen determinados por los ceros positivos de la función de Bessel de orden cuatro,

\[
J_4(j_{4,m})=0,
\]

y todos los modos decaen exponencialmente con el tiempo, como corresponde a la ecuación del calor.











\section*{Problema 9 (Ecuación de ondas en 2D)}
Considere la ecuación de ondas en dos dimensiones espaciales:
\[
\frac{\partial^2 u}{\partial t^2} = c^2 \left( \frac{\partial^2 u}{\partial x^2} + \frac{\partial^2 u}{\partial y^2} \right),
\]
definida en el rectángulo \(0 < x < a\), \(0 < y < b\), con condiciones de contorno homogéneas:
\[
u(0, y, t) = u(a, y, t) = 0, \qquad u(x, 0, t) = u(x, b, t) = 0,
\]
y condiciones iniciales:
\[
u(x, y, 0) = f(x, y), \qquad \frac{\partial u}{\partial t}(x, y, 0) = g(x, y).
\]
\begin{enumerate}
    \item[(a)] Proponga una solución por separación de variables \(u(x,y,t) = X(x)Y(y)T(t)\) y obtenga las ecuaciones diferenciales ordinarias para \(X\), \(Y\) y \(T\).
    \item[(b)] Determine las autofunciones espaciales \(X_m(x)\) y \(Y_n(y)\) y los correspondientes autovalores.
    \item[(c)] Escriba la solución general como una serie doble.
    \item[(d)] En el caso particular \(f(x,y) = \sen\left(\frac{2\pi x}{a}\right) \sen\left(\frac{3\pi y}{b}\right)\) y \(g(x,y) = 0\), obtenga la solución explícita.
\end{enumerate}

\vspace{0.5cm}
{\bf Soluci\'on }{\em 9)}
\vspace{0.5cm}

Consideremos la ecuación de ondas en dos dimensiones

\begin{equation}
\frac{\partial^2u}{\partial t^2}
=
c^2
\left(
\frac{\partial^2u}{\partial x^2}
+
\frac{\partial^2u}{\partial y^2}
\right),
\end{equation}

definida en el rectángulo

\[
0<x<a,\qquad
0<y<b,
\]

con condiciones de contorno

\begin{equation}
u(0,y,t)=u(a,y,t)=0,
\end{equation}

\begin{equation}
u(x,0,t)=u(x,b,t)=0,
\end{equation}

y condiciones iniciales

\begin{equation}
u(x,y,0)=f(x,y),
\end{equation}

\begin{equation}
\frac{\partial u}{\partial t}(x,y,0)=g(x,y).
\end{equation}

\subsubsection*{(a) Separación de variables}

Buscamos una solución de la forma

\begin{equation}
u(x,y,t)=X(x)Y(y)T(t).
\end{equation}

Sustituyendo en la ecuación de ondas,

\[
XYT''
=
c^2
\left(
X''YT
+
XY''T
\right).
\]

Dividiendo por $XYT$,

\begin{equation}
\frac{T''}{c^2T}
=
\frac{X''}{X}
+
\frac{Y''}{Y}.
\end{equation}

El miembro izquierdo depende únicamente del tiempo, mientras que el derecho depende únicamente de las coordenadas espaciales.

Introducimos una constante de separación $-\lambda$,

\begin{equation}
\frac{T''}{c^2T}
=
-\lambda,
\end{equation}

obteniéndose

\begin{equation}
\frac{X''}{X}
+
\frac{Y''}{Y}
=
-\lambda.
\end{equation}

Separamos nuevamente mediante otra constante $-\mu$,

\begin{equation}
\frac{X''}{X}
=
-\mu,
\end{equation}

y

\begin{equation}
\frac{Y''}{Y}
=
-(\lambda-\mu).
\end{equation}

Por consiguiente,

\[
\boxed{
X''+\mu X=0,
}
\]

\[
\boxed{
Y''+(\lambda-\mu)Y=0,
}
\]

y

\[
\boxed{
T''+c^2\lambda T=0.
}
\]

Éstas son las ecuaciones diferenciales buscadas.

\subsubsection*{(b) Autofunciones espaciales}

\paragraph{Dirección $x$.}

Debemos resolver

\[
X''+\mu X=0,
\]

con

\[
X(0)=X(a)=0.
\]

La solución general es

\[
X(x)=A\cos(\sqrt{\mu}\,x)
+
B\sin(\sqrt{\mu}\,x).
\]

La condición en $x=0$ implica

\[
A=0.
\]

La condición en $x=a$ exige

\[
\sin(\sqrt{\mu}\,a)=0,
\]

de donde

\begin{equation}
\sqrt{\mu}
=
\frac{m\pi}{a},
\qquad
m=1,2,\ldots
\end{equation}

Así,

\begin{equation}
X_m(x)
=
\sin\!\left(\frac{m\pi x}{a}\right).
\end{equation}

\paragraph{Dirección $y$.}

De manera completamente análoga,

\[
Y''+\nu Y=0,
\]

con

\[
Y(0)=Y(b)=0,
\]

produce

\begin{equation}
\nu
=
\left(\frac{n\pi}{b}\right)^2,
\qquad
n=1,2,\ldots
\end{equation}

y

\begin{equation}
Y_n(y)
=
\sin\!\left(\frac{n\pi y}{b}\right).
\end{equation}

Como

\[
\lambda=\mu+\nu,
\]

resulta

\begin{equation}
\lambda_{mn}
=
\left(\frac{m\pi}{a}\right)^2
+
\left(\frac{n\pi}{b}\right)^2.
\end{equation}

\subsubsection*{Ecuación temporal}

Cada modo espacial satisface

\begin{equation}
T_{mn}''
+
\omega_{mn}^2T_{mn}=0,
\end{equation}

donde

\begin{equation}
\omega_{mn}
=
c
\sqrt{
\left(\frac{m\pi}{a}\right)^2
+
\left(\frac{n\pi}{b}\right)^2
}.
\end{equation}

La solución general es

\begin{equation}
T_{mn}(t)
=
A_{mn}\cos(\omega_{mn}t)
+
B_{mn}\sin(\omega_{mn}t).
\end{equation}

\subsubsection*{(c) Solución general}

Superponiendo todos los modos normales,

\begin{equation}
\boxed{
u(x,y,t)
=
\sum_{m=1}^{\infty}
\sum_{n=1}^{\infty}
\left[
A_{mn}\cos(\omega_{mn}t)
+
B_{mn}\sin(\omega_{mn}t)
\right]
\sin\!\left(\frac{m\pi x}{a}\right)
\sin\!\left(\frac{n\pi y}{b}\right).
}
\end{equation}

Los coeficientes se obtienen a partir de las condiciones iniciales.

De

\[
u(x,y,0)=f(x,y),
\]

se obtiene

\begin{equation}
A_{mn}
=
\frac{4}{ab}
\int_0^a
\int_0^b
f(x,y)
\sin\!\left(\frac{m\pi x}{a}\right)
\sin\!\left(\frac{n\pi y}{b}\right)
dx\,dy.
\end{equation}

Derivando respecto del tiempo,

\[
u_t(x,y,t)
=
\sum_{m,n}
\omega_{mn}
B_{mn}
\cos(\omega_{mn}t)
\sin\!\left(\frac{m\pi x}{a}\right)
\sin\!\left(\frac{n\pi y}{b}\right)
+\cdots
\]

y evaluando en $t=0$,

\[
u_t(x,y,0)=g(x,y),
\]

obtenemos

\begin{equation}
B_{mn}
=
\frac{4}
{ab\,\omega_{mn}}
\int_0^a
\int_0^b
g(x,y)
\sin\!\left(\frac{m\pi x}{a}\right)
\sin\!\left(\frac{n\pi y}{b}\right)
dx\,dy.
\end{equation}

\subsubsection*{(d) Caso particular}

Supongamos ahora

\begin{equation}
f(x,y)
=
\sin\!\left(\frac{2\pi x}{a}\right)
\sin\!\left(\frac{3\pi y}{b}\right),
\end{equation}

y

\[
g(x,y)=0.
\]

Por ortogonalidad de las funciones seno,

\[
A_{mn}=0,
\qquad
(m,n)\neq(2,3),
\]

mientras que

\[
A_{23}=1,
\]

y todos los coeficientes

\[
B_{mn}=0.
\]

Por lo tanto, únicamente sobrevive un modo normal.

Su frecuencia angular es

\begin{equation}
\omega_{23}
=
c
\sqrt{
\left(\frac{2\pi}{a}\right)^2
+
\left(\frac{3\pi}{b}\right)^2
}.
\end{equation}

Finalmente,

\begin{equation}
\boxed{
u(x,y,t)
=
\cos(\omega_{23}t)
\sin\!\left(\frac{2\pi x}{a}\right)
\sin\!\left(\frac{3\pi y}{b}\right),
}
\end{equation}

donde

\begin{equation}
\boxed{
\omega_{23}
=
c
\sqrt{
\left(\frac{2\pi}{a}\right)^2
+
\left(\frac{3\pi}{b}\right)^2
}.
}
\end{equation}

Esta solución satisface tanto la ecuación de ondas como las condiciones de contorno homogéneas y las condiciones iniciales especificadas.








\section*{Problema 10 (Ondas en una membrana circular)}
Una membrana circular de radio \(R\) se encuentra fija en su borde. Su desplazamiento vertical \(u(r, \theta, t)\) satisface la ecuación de ondas en coordenadas polares:
\[
\frac{\partial^2 u}{\partial t^2} = c^2 \left( \frac{1}{r} \frac{\partial}{\partial r} \left( r \frac{\partial u}{\partial r} \right) + \frac{1}{r^2} \frac{\partial^2 u}{\partial \theta^2} \right),
\]
con condición de borde \(u(R, \theta, t) = 0\).
\begin{enumerate}
    \item[(a)] Proponga una solución de la forma \(u(r,\theta,t) = R(r)\Theta(\theta)T(t)\) y obtenga las ecuaciones diferenciales para cada función.
    \item[(b)] Resuelva la ecuación para \(\Theta(\theta)\) imponiendo periodicidad en \(\theta\).
    \item[(c)] Muestre que la ecuación radial se reduce a la ecuación de Bessel y escriba su solución general (en términos de funciones de Bessel).
    \item[(d)] Aplique la condición de borde para determinar los autovalores permitidos y escriba la forma general de la solución.
\end{enumerate}

\vspace{0.5cm}
{\bf Soluci\'on }{\em 10)}
\vspace{0.5cm}

Consideremos la ecuación de ondas en coordenadas polares

\begin{equation}
\frac{\partial^2u}{\partial t^2}
=
c^2
\left[
\frac{1}{r}\frac{\partial}{\partial r}
\left(
r\frac{\partial u}{\partial r}
\right)
+
\frac{1}{r^2}
\frac{\partial^2u}{\partial\theta^2}
\right],
\end{equation}

definida sobre un disco de radio $R$, con condición de borde

\begin{equation}
u(R,\theta,t)=0.
\end{equation}

Se supone además que la solución es periódica en la variable angular,

\[
u(r,\theta+2\pi,t)=u(r,\theta,t),
\]

y permanece finita en el origen.

\subsubsection*{(a) Separación de variables}

Buscamos una solución de la forma

\begin{equation}
u(r,\theta,t)=R(r)\Theta(\theta)T(t).
\end{equation}

Sustituyendo en la ecuación de ondas,

\[
R\Theta T''
=
c^2
\left[
\frac1r
\frac{d}{dr}
\left(
rR'
\right)
\Theta T
+
\frac1{r^2}
R\Theta''T
\right].
\]

Dividiendo por $c^2R\Theta T$,

\begin{equation}
\frac{T''}{c^2T}
=
\frac1R
\left(
R''+\frac1rR'
\right)
+
\frac1{r^2}
\frac{\Theta''}{\Theta}.
\end{equation}

El miembro izquierdo depende únicamente de $t$, mientras que el derecho depende sólo de las variables espaciales.

Introducimos una constante de separación $-\lambda$,

\begin{equation}
\frac{T''}{c^2T}
=
-\lambda,
\end{equation}

obteniéndose

\begin{equation}
\frac1R
\left(
R''+\frac1rR'
\right)
+
\frac1{r^2}
\frac{\Theta''}{\Theta}
=
-\lambda.
\end{equation}

Multiplicando por $r^2$,

\[
r^2
\frac{R''+\frac1rR'}{R}
+
\frac{\Theta''}{\Theta}
=
-\lambda r^2.
\]

Separamos nuevamente mediante una constante $-m^2$,

\begin{equation}
\frac{\Theta''}{\Theta}
=
-m^2,
\end{equation}

de donde obtenemos las ecuaciones

\begin{equation}
\boxed{
\Theta''+m^2\Theta=0,
}
\end{equation}

\begin{equation}
\boxed{
r^2R''
+rR'
+
(\lambda r^2-m^2)R=0,
}
\end{equation}

y

\begin{equation}
\boxed{
T''+c^2\lambda T=0.
}
\end{equation}

\subsubsection*{(b) Solución de la ecuación angular}

La ecuación angular es

\[
\Theta''+m^2\Theta=0.
\]

Su solución general es

\begin{equation}
\Theta(\theta)
=
A\cos(m\theta)
+
B\sin(m\theta).
\end{equation}

La periodicidad

\[
\Theta(\theta+2\pi)=\Theta(\theta)
\]

exige que

\[
m=0,\pm1,\pm2,\ldots
\]

Por lo tanto,

\begin{equation}
\boxed{
m\in\mathbb Z,
}
\end{equation}

y las autofunciones angulares son

\begin{equation}
\boxed{
\Theta_m(\theta)
=
A_m\cos(m\theta)
+
B_m\sin(m\theta).
}
\end{equation}

\subsubsection*{(c) Ecuación radial}

La ecuación radial es

\begin{equation}
r^2R''
+rR'
+
(\lambda r^2-m^2)R=0.
\end{equation}

Introducimos la variable

\[
x=\sqrt{\lambda}\,r.
\]

Como

\[
\frac{dR}{dr}
=
\sqrt{\lambda}
\frac{dR}{dx},
\qquad
\frac{d^2R}{dr^2}
=
\lambda
\frac{d^2R}{dx^2},
\]

la ecuación anterior se transforma en

\begin{equation}
x^2R''
+xR'
+
(x^2-m^2)R=0.
\end{equation}

Ésta es precisamente la ecuación de Bessel de orden $m$.

Su solución general es

\begin{equation}
\boxed{
R(r)
=
A
J_m(\sqrt{\lambda}\,r)
+
B
Y_m(\sqrt{\lambda}\,r),
}
\end{equation}

donde $J_m$ y $Y_m$ son las funciones de Bessel de primera y segunda especie.

Como

\[
Y_m(r)\rightarrow\infty,
\qquad
r\rightarrow0,
\]

y la membrana debe tener desplazamiento finito en el centro, necesariamente

\[
B=0.
\]

Por consiguiente,

\begin{equation}
\boxed{
R(r)
=
J_m(\sqrt{\lambda}\,r).
}
\end{equation}

\subsubsection*{(d) Autovalores y solución general}

La condición de borde

\[
u(R,\theta,t)=0
\]

implica

\[
R(R)=0,
\]

es decir,

\begin{equation}
J_m(\sqrt{\lambda}\,R)=0.
\end{equation}

Denotemos por

\[
j_{m,n},
\qquad
n=1,2,\ldots,
\]

el $n$-ésimo cero positivo de la función de Bessel $J_m$.

Entonces

\begin{equation}
\sqrt{\lambda_{mn}}
=
\frac{j_{m,n}}{R},
\end{equation}

o equivalentemente,

\begin{equation}
\boxed{
\lambda_{mn}
=
\left(
\frac{j_{m,n}}{R}
\right)^2.
}
\end{equation}

La ecuación temporal queda

\begin{equation}
T_{mn}''
+\omega_{mn}^2T_{mn}=0,
\end{equation}

con

\begin{equation}
\omega_{mn}
=
c
\frac{j_{m,n}}{R}.
\end{equation}

Su solución general es

\begin{equation}
T_{mn}(t)
=
C_{mn}\cos(\omega_{mn}t)
+
D_{mn}\sin(\omega_{mn}t).
\end{equation}

Finalmente, la solución general del problema se obtiene superponiendo todos los modos normales,

\begin{equation}
\boxed{
\begin{aligned}
u(r,\theta,t)
=
\sum_{m=0}^{\infty}
\sum_{n=1}^{\infty}
&
J_m
\!\left(
\frac{j_{m,n}}{R}r
\right)
\Big[
A_{mn}\cos(m\theta)
+
B_{mn}\sin(m\theta)
\Big]
\\
&\times
\left[
C_{mn}\cos(\omega_{mn}t)
+
D_{mn}\sin(\omega_{mn}t)
\right],
\end{aligned}
}
\end{equation}

donde

\begin{equation}
\boxed{
\omega_{mn}
=
c
\frac{j_{m,n}}{R}.
}
\end{equation}

Los coeficientes $A_{mn}$, $B_{mn}$, $C_{mn}$ y $D_{mn}$ quedan completamente determinados por las condiciones iniciales de desplazamiento y velocidad de la membrana.

\bigskip

\noindent
\textbf{Respuesta final.}

\begin{itemize}
\item La ecuación angular conduce a los modos
\[
\Theta_m(\theta)=A\cos(m\theta)+B\sin(m\theta),
\qquad
m\in\mathbb Z.
\]

\item La ecuación radial es la ecuación de Bessel

\[
r^2R''
+rR'
+
(\lambda r^2-m^2)R=0,
\]

cuya solución regular es

\[
R(r)=J_m(\sqrt{\lambda}\,r).
\]

\item Los autovalores vienen dados por los ceros de las funciones de Bessel,

\[
J_m(\sqrt{\lambda}\,R)=0,
\]

es decir,

\[
\lambda_{mn}
=
\left(
\frac{j_{m,n}}{R}
\right)^2.
\]

\item La solución general es una superposición de modos normales de la forma

\[
J_m\!\left(\frac{j_{m,n}}{R}r\right)
\left[
A\cos(m\theta)+B\sin(m\theta)
\right]
\left[
C\cos(\omega_{mn}t)+D\sin(\omega_{mn}t)
\right].
\]
\end{itemize}







\section*{Problema 11 (Ondas en un sector circular)}
Considere la ecuación de ondas en un sector circular de radio \(R\) y ángulo \(\alpha\) (\(0 < \alpha < 2\pi\)):
\[
\frac{\partial^2 u}{\partial t^2} = c^2 \left( \frac{1}{r} \frac{\partial}{\partial r} \left( r \frac{\partial u}{\partial r} \right) + \frac{1}{r^2} \frac{\partial^2 u}{\partial \theta^2} \right),
\]
con condiciones de contorno:
\[
u(r, 0, t) = 0, \qquad u(r, \alpha, t) = 0, \qquad u(R, \theta, t) = 0,
\]
y condiciones iniciales \(u(r, \theta, 0) = f(r, \theta)\), \(\partial_t u(r, \theta, 0) = g(r, \theta)\).
\begin{enumerate}
    \item[(a)] Proponga una solución por separación de variables y obtenga las ecuaciones para \(R(r)\), \(\Theta(\theta)\) y \(T(t)\).
    \item[(b)] Resuelva la ecuación angular y determine los autovalores \(\lambda_n\).
    \item[(c)] Muestre que la ecuación radial se reduce a la ecuación de Bessel y determine los autovalores radiales a partir de la condición de borde en \(r = R\).
    \item[(d)] Escriba la solución general como una serie doble.
\end{enumerate}

\vspace{0.5cm}
{\bf Soluci\'on }{\em 11)}
\vspace{0.5cm}

Consideremos la ecuación de ondas en coordenadas polares

\begin{equation}
\frac{\partial^2u}{\partial t^2}
=
c^2
\left[
\frac{1}{r}\frac{\partial}{\partial r}
\left(
r\frac{\partial u}{\partial r}
\right)
+
\frac{1}{r^2}
\frac{\partial^2u}{\partial\theta^2}
\right],
\end{equation}

definida en un sector circular

\[
0<r<R,
\qquad
0<\theta<\alpha,
\]

con condiciones de contorno

\begin{equation}
u(r,0,t)=0,
\qquad
u(r,\alpha,t)=0,
\qquad
u(R,\theta,t)=0,
\end{equation}

y condiciones iniciales

\begin{equation}
u(r,\theta,0)=f(r,\theta),
\qquad
\frac{\partial u}{\partial t}(r,\theta,0)=g(r,\theta).
\end{equation}

\subsubsection*{(a) Separación de variables}

Buscamos una solución de la forma

\begin{equation}
u(r,\theta,t)=R(r)\Theta(\theta)T(t).
\end{equation}

Sustituyendo en la ecuación de ondas,

\[
R\Theta T''
=
c^2
\left[
\frac1r
\frac{d}{dr}(rR')
\Theta T
+
\frac1{r^2}
R\Theta''T
\right].
\]

Dividiendo por $c^2R\Theta T$,

\begin{equation}
\frac{T''}{c^2T}
=
\frac1R
\left(
R''+\frac1rR'
\right)
+
\frac1{r^2}
\frac{\Theta''}{\Theta}.
\end{equation}

El miembro izquierdo depende únicamente del tiempo, mientras que el derecho depende únicamente de las variables espaciales.

Introducimos una constante de separación $-\mu^2$,

\begin{equation}
\frac{T''}{c^2T}
=
-\mu^2,
\end{equation}

obteniéndose

\begin{equation}
\frac1R
\left(
R''+\frac1rR'
\right)
+
\frac1{r^2}
\frac{\Theta''}{\Theta}
=
-\mu^2.
\end{equation}

Multiplicando por $r^2$,

\[
r^2
\frac{R''+\frac1rR'}{R}
+
\frac{\Theta''}{\Theta}
=
-\mu^2r^2.
\]

Separamos nuevamente mediante una constante $-\lambda$,

\begin{equation}
\frac{\Theta''}{\Theta}
=
-\lambda,
\end{equation}

obteniéndose las ecuaciones

\begin{equation}
\boxed{
\Theta''+\lambda\Theta=0,
}
\end{equation}

\begin{equation}
\boxed{
r^2R''
+rR'
+
(\mu^2r^2-\lambda)R=0,
}
\end{equation}

y

\begin{equation}
\boxed{
T''+c^2\mu^2T=0.
}
\end{equation}

\subsubsection*{(b) Ecuación angular}

Debemos resolver

\begin{equation}
\Theta''+\lambda\Theta=0,
\end{equation}

con

\[
\Theta(0)=0,
\qquad
\Theta(\alpha)=0.
\]

La solución general es

\[
\Theta(\theta)
=
A\cos(\sqrt{\lambda}\,\theta)
+
B\sin(\sqrt{\lambda}\,\theta).
\]

La condición en $\theta=0$ implica

\[
A=0,
\]

de modo que

\[
\Theta(\theta)
=
B
\sin(\sqrt{\lambda}\,\theta).
\]

La condición en $\theta=\alpha$ exige

\[
\sin(\sqrt{\lambda}\,\alpha)=0,
\]

por lo que

\begin{equation}
\sqrt{\lambda}
=
\frac{n\pi}{\alpha},
\qquad
n=1,2,\ldots
\end{equation}

En consecuencia,

\begin{equation}
\boxed{
\lambda_n
=
\left(
\frac{n\pi}{\alpha}
\right)^2,
}
\end{equation}

y las autofunciones angulares son

\begin{equation}
\boxed{
\Theta_n(\theta)
=
\sin
\left(
\frac{n\pi\theta}{\alpha}
\right).
}
\end{equation}

\subsubsection*{(c) Ecuación radial}

La ecuación radial queda

\begin{equation}
r^2R''
+rR'
+
\left[
\mu^2r^2
-
\left(
\frac{n\pi}{\alpha}
\right)^2
\right]
R=0.
\end{equation}

Introduciendo la variable

\[
x=\mu r,
\]

la ecuación se transforma en

\begin{equation}
x^2R''
+xR'
+
\left(
x^2-\nu^2
\right)
R=0,
\end{equation}

donde

\[
\nu=\frac{n\pi}{\alpha}.
\]

Ésta es la ecuación de Bessel de orden $\nu$.

La solución general es

\begin{equation}
R(r)
=
AJ_\nu(\mu r)
+
BY_\nu(\mu r).
\end{equation}

Como la solución debe permanecer finita en el origen,

\[
Y_\nu(r)\rightarrow\infty,
\qquad
r\rightarrow0,
\]

debemos tomar

\[
B=0.
\]

Por lo tanto,

\begin{equation}
\boxed{
R(r)
=
J_\nu(\mu r).
}
\end{equation}

La condición de borde

\[
u(R,\theta,t)=0
\]

implica

\begin{equation}
J_\nu(\mu R)=0.
\end{equation}

Si denotamos por

\[
j_{\nu,m},
\qquad
m=1,2,\ldots,
\]

el $m$-ésimo cero positivo de la función de Bessel de orden $\nu$, entonces

\begin{equation}
\boxed{
\mu_{nm}
=
\frac{j_{\nu,m}}{R},
}
\end{equation}

con

\[
\nu=\frac{n\pi}{\alpha}.
\]

Los autovalores radiales quedan completamente determinados por los ceros de las funciones de Bessel.

\subsubsection*{Parte temporal}

La ecuación temporal es

\begin{equation}
T''+\omega_{nm}^2T=0,
\end{equation}

donde

\begin{equation}
\boxed{
\omega_{nm}
=
c
\frac{j_{\nu,m}}{R}.
}
\end{equation}

Su solución general es

\begin{equation}
T_{nm}(t)
=
C_{nm}\cos(\omega_{nm}t)
+
D_{nm}\sin(\omega_{nm}t).
\end{equation}

\subsubsection*{(d) Solución general}

Combinando los resultados anteriores, obtenemos

\begin{equation}
\boxed{
\begin{aligned}
u(r,\theta,t)
=
\sum_{n=1}^{\infty}
\sum_{m=1}^{\infty}
&
\left[
A_{nm}\cos(\omega_{nm}t)
+
B_{nm}\sin(\omega_{nm}t)
\right]
\\
&\times
J_{\frac{n\pi}{\alpha}}
\left(
\frac{j_{\frac{n\pi}{\alpha},m}}{R}r
\right)
\sin
\left(
\frac{n\pi\theta}{\alpha}
\right),
\end{aligned}
}
\end{equation}

donde

\begin{equation}
\omega_{nm}
=
c
\frac{j_{\frac{n\pi}{\alpha},m}}{R}.
\end{equation}

Los coeficientes $A_{nm}$ y $B_{nm}$ se determinan expandiendo las condiciones iniciales $f(r,\theta)$ y $g(r,\theta)$ en la base ortogonal formada por las funciones

\[
J_{\frac{n\pi}{\alpha}}
\left(
\frac{j_{\frac{n\pi}{\alpha},m}}{R}r
\right)
\sin
\left(
\frac{n\pi\theta}{\alpha}
\right),
\]

utilizando la ortogonalidad de las funciones seno en el intervalo $(0,\alpha)$ y de las funciones de Bessel en el intervalo $(0,R)$.











\section*{Problema 12 (Membrana circular con dos densidades)}
Una membrana circular de radio \(R\) está compuesta por dos materiales de distintas densidades. La densidad superficial \(\rho(r)\) es:
\[
\rho(r) = 
\begin{cases}
\rho_1, & 0 < r < a, \\
\rho_2, & a < r < R,
\end{cases}
\]
con \(\rho_1 \neq \rho_2\). La ecuación de ondas para cada región es:
\[
\frac{\partial^2 u_i}{\partial t^2} = c_i^2 \nabla^2 u_i, \qquad i = 1, 2,
\]
donde \(c_i^2 = T/\rho_i\) (siendo \(T\) la tensión uniforme). La membrana está fija en su borde: \(u_2(R, \theta, t) = 0\). Además, en la interfaz \(r = a\) se deben cumplir las condiciones de continuidad:
\[
u_1(a, \theta, t) = u_2(a, \theta, t), \qquad \frac{\partial u_1}{\partial r}(a, \theta, t) = \frac{\partial u_2}{\partial r}(a, \theta, t).
\]
Suponga soluciones de la forma \(u_i(r, \theta, t) = R_i(r) e^{i m \theta} e^{i \omega t}\) (con \(m \in \mathbb{Z}\)).
\begin{enumerate}
    \item[(a)] Escriba la ecuación diferencial para \(R_i(r)\) en cada región y muestre que se trata de la ecuación de Bessel.
    \item[(b)] Proponga la forma de la solución radial en cada región (use funciones de Bessel y Neumann, o Hankel, según convenga) y justifique la elección en el centro.
    \item[(c)] Aplique las condiciones de contorno e interfaz para obtener un sistema de ecuaciones para los coeficientes.
    \item[(d)] Escriba la ecuación que determina las frecuencias propias \(\omega\) (no es necesario resolverla explícitamente, solo plantear el determinante).
\end{enumerate}

\vspace{0.5cm}
{\bf Soluci\'on }{\em 12)}
\vspace{0.5cm}

La membrana circular está formada por dos materiales con distintas densidades superficiales,

\[
\rho(r)=
\begin{cases}
\rho_1, & 0<r<a,\\
\rho_2, & a<r<R,
\end{cases}
\]

por lo que la velocidad de propagación es distinta en cada región,

\[
c_i^2=\frac{T}{\rho_i},
\qquad
i=1,2.
\]

En consecuencia, la ecuación de ondas debe resolverse independientemente en cada dominio,

\begin{equation}
\frac{\partial^2u_i}{\partial t^2}
=
c_i^2\nabla^2u_i,
\qquad
i=1,2.
\end{equation}

Se supone una solución de la forma

\begin{equation}
u_i(r,\theta,t)
=
R_i(r)e^{im\theta}e^{i\omega t},
\qquad
m\in\mathbb Z.
\end{equation}

Recordemos que el Laplaciano en coordenadas polares es

\begin{equation}
\nabla^2u
=
\frac1r
\frac{\partial}{\partial r}
\left(
r\frac{\partial u}{\partial r}
\right)
+
\frac1{r^2}
\frac{\partial^2u}{\partial\theta^2}.
\end{equation}

\subsubsection*{(a) Ecuación radial}

Sustituyendo la forma propuesta en la ecuación de ondas,

\[
-\omega^2R_ie^{im\theta}
=
c_i^2
\left[
\frac1r
\frac{d}{dr}
\left(
rR_i'
\right)
-
\frac{m^2}{r^2}R_i
\right]
e^{im\theta}.
\]

Dividiendo por $e^{im\theta}$ obtenemos

\[
\frac1r
\frac{d}{dr}
\left(
rR_i'
\right)
+
\left(
\frac{\omega^2}{c_i^2}
-
\frac{m^2}{r^2}
\right)
R_i
=0.
\]

Definiendo

\[
k_i=\frac{\omega}{c_i},
\]

la ecuación radial queda

\begin{equation}
R_i''
+\frac1rR_i'
+
\left(
k_i^2-\frac{m^2}{r^2}
\right)
R_i=0.
\end{equation}

Multiplicando por $r^2$,

\begin{equation}
\boxed{
r^2R_i''
+rR_i'
+
(k_i^2r^2-m^2)R_i=0.
}
\end{equation}

Ésta es la ecuación de Bessel de orden $m$.

\subsubsection*{(b) Soluciones radiales}

La solución general de la ecuación de Bessel es

\begin{equation}
R_i(r)
=
A_iJ_m(k_ir)
+
B_iY_m(k_ir),
\end{equation}

donde $J_m$ y $Y_m$ son las funciones de Bessel de primera y segunda especie.

\paragraph{Región interior ($0<r<a$).}

Como

\[
Y_m(r)\longrightarrow\infty,
\qquad
r\rightarrow0,
\]

la solución debe permanecer finita en el centro de la membrana. Por ello,

\[
B_1=0,
\]

y

\begin{equation}
\boxed{
R_1(r)=A\,J_m(k_1r).
}
\end{equation}

\paragraph{Región exterior ($a<r<R$).}

En esta región el origen no pertenece al dominio, por lo que ambas soluciones son admisibles. Así,

\begin{equation}
\boxed{
R_2(r)
=
C\,J_m(k_2r)
+
D\,Y_m(k_2r).
}
\end{equation}

\subsubsection*{(c) Condiciones de contorno e interfaz}

\paragraph{Condición en el borde exterior.}

La membrana está fija en $r=R$,

\[
u_2(R,\theta,t)=0,
\]

lo que implica

\begin{equation}
\boxed{
CJ_m(k_2R)
+
DY_m(k_2R)
=0.
}
\end{equation}

\paragraph{Continuidad del desplazamiento.}

En la interfaz $r=a$,

\[
u_1(a,\theta,t)=u_2(a,\theta,t),
\]

de donde

\begin{equation}
\boxed{
AJ_m(k_1a)
=
CJ_m(k_2a)
+
DY_m(k_2a).
}
\end{equation}

\paragraph{Continuidad de la derivada radial.}

La segunda condición es

\[
\frac{\partial u_1}{\partial r}
=
\frac{\partial u_2}{\partial r},
\qquad
r=a.
\]

Como

\[
\frac{d}{dr}
J_m(kr)
=
kJ_m'(kr),
\]

se obtiene

\begin{equation}
\boxed{
Ak_1J_m'(k_1a)
=
Ck_2J_m'(k_2a)
+
Dk_2Y_m'(k_2a).
}
\end{equation}

Las tres condiciones anteriores constituyen un sistema homogéneo para los coeficientes $A$, $C$ y $D$,

\begin{equation}
\begin{pmatrix}
J_m(k_1a)
&
-J_m(k_2a)
&
-Y_m(k_2a)
\\[2mm]
k_1J_m'(k_1a)
&
-k_2J_m'(k_2a)
&
-k_2Y_m'(k_2a)
\\[2mm]
0
&
J_m(k_2R)
&
Y_m(k_2R)
\end{pmatrix}
\begin{pmatrix}
A\\
C\\
D
\end{pmatrix}
=
\begin{pmatrix}
0\\
0\\
0
\end{pmatrix}.
\end{equation}

\subsubsection*{(d) Frecuencias propias}

El sistema anterior posee soluciones no triviales únicamente cuando el determinante de la matriz de coeficientes se anula. En consecuencia, las frecuencias propias de la membrana quedan determinadas por la ecuación trascendente

\begin{equation}
\boxed{
\begin{vmatrix}
J_m(k_1a)
&
-J_m(k_2a)
&
-Y_m(k_2a)
\\[2mm]
k_1J_m'(k_1a)
&
-k_2J_m'(k_2a)
&
-k_2Y_m'(k_2a)
\\[2mm]
0
&
J_m(k_2R)
&
Y_m(k_2R)
\end{vmatrix}
=0,
}
\end{equation}

donde

\[
k_1=\frac{\omega}{c_1},
\qquad
k_2=\frac{\omega}{c_2}.
\]

Esta ecuación determina implícitamente las frecuencias propias $\omega$ del sistema. En general, no admite solución analítica y debe resolverse mediante métodos numéricos.

\bigskip

\noindent
\textbf{Observación.} En el caso particular $\rho_1=\rho_2$, se tiene $c_1=c_2$, desaparece la interfaz y la ecuación característica se reduce a la condición habitual

\[
J_m(kR)=0,
\]

que corresponde a la membrana circular homogénea.





\end{document}